\subsection{Auxiliary classical identities}
The rational-letter terms of the preceding coefficient list do not
require separate GPL evaluations. Put $L=\ln(1+X)$, $M=\ln X$ and
$v=1/(1+X)$, with $X>0$. The elementary reductions are
\begin{align}
 G_{-1}&=L,\qquad G_0=M,\qquad
 G_{-1,-1}=\tfrac12L^2,\qquad G_{-1,-1,-1}=\tfrac16L^3,\\
 G_{0,-1}&=-\operatorname{Li}_2(-X),\qquad
 G_{-1,0}=LM+\operatorname{Li}_2(-X),\\
 G_{1,0}&=\frac{\pi^2}{6}-\operatorname{Li}_2(1-X),\\
 G_{0,-1,0}&=-M\operatorname{Li}_2(-X)+2\operatorname{Li}_3(-X),\\
 G_{0,1,0}&=-M\mathcal L_2(X)+2\mathcal L_3(X).
\end{align}
Here $\mathcal L_n(X)=\operatorname{Re}\operatorname{Li}_n(X+\ii0)$.
The real boundary value is fixed by integrating
$\dd\mathcal L_2/\dd X=-\ln|1-X|/X$ and
$\dd\mathcal L_3/\dd X=\mathcal L_2/X$ from the interval $0<X<1$.
It is not selected by fitting numerical values. The combination
for $G_{1,0}$ is regular at $X=1$.

One additional classical combination suffices for the three words
with two letters $-1$ and one zero:
\begin{align}
 T(X)&=\tfrac12L^2M-\tfrac13L^3
       -L\operatorname{Li}_2(v)-\operatorname{Li}_3(v)+\zeta(3),\\
 G_{0,-1,-1}&=T(X),\\
 G_{-1,0,-1}&=-L\operatorname{Li}_2(-X)-2T(X),\\
 G_{-1,-1,0}&=\tfrac12L^2M+L\operatorname{Li}_2(-X)+T(X).
\end{align}
For example, substituting $v=1/(1+X)$ into
$\tfrac12\int_0^X\ln^2(1+t)\,\dd t/t$ gives the displayed $T$.
Its derivative is $L^2/(2X)$ and its lower limit is zero.
The last two lines then follow from the shuffle product.
All dilogarithms and trilogarithms with argument $v$ are on
$0<v<1$. These substitutions reduce the number of nonclassical
words in the partial sum without changing any rational coefficient
or the photon-mass convention.

The repeated square-root kernel can also be integrated classically.
Use the same $q$ as in \eqref{eq:comp-ir-short-function}. Since
$\omega_r=-\dd q/q$, integrating that independently derived expression
once more gives
\begin{equation}
 G_{r,r,-1}=\operatorname{Li}_3(-q)
 -\frac19\operatorname{Li}_3(-q^3)-\frac16\ln^3q
 +\frac{\pi^2}{18}\ln q+\frac23\zeta(3).
\end{equation}
Differentiating with respect to $q$ gives
$-\mathcal G(X(q))/q$. At $q=1$ the expression vanishes, using
$\operatorname{Li}_3(-1)=-3\zeta(3)/4$, which fixes the constant.
Its two trilogarithms have negative real arguments.
The remaining rational-letter word also has a short classical form.
With $u=1+X$, its inner weight-two function is
$\operatorname{Li}_2(-u)+\ln u\ln(1+u)+\pi^2/12$.
Integrating it with $\dd u/u$ gives
\begin{equation}
 G_{-1,-2,-1}=2\operatorname{Li}_3(-1-X)
 -L\operatorname{Li}_2(-1-X)+\frac{\pi^2}{12}L
 +\frac32\zeta(3).
\end{equation}
The last constant follows from the zero value at $X=0$.
Together with \eqref{eq:comp-ir-short-function}, these reductions
leave only six nonclassical words in this partial sum:
\begin{equation}
 \left\{G_{-4,r,-1},G_{0,r,-1},G_{r,-1,-1},
 G_{r,-1,0},G_{r,-2,-1},G_{r,0,-1}\right\}.
\end{equation}

\paragraph{A classical primitive for the squared logarithm.}
The word $G_{r,-1,-1}$ can be removed from this list without
introducing another iterated integral. Define
\begin{align}
 a&=e^{\ii\pi/3},\quad b=\overline a,\quad X=(1-q)^2/q,\quad 0<q<1,\\
 A&=1-aq,\quad B=1-bq,\quad R=A/B.
\end{align}
The quadratic factorization gives $AB=1-q+q^2$ and hence
$\ln(1+X)=\log A+\log B-\ln q$.
Here and below $\log$ is the principal complex logarithm.
Introduce the classical function
\begin{equation}
 \mathcal J(z,\rho)=\log^2z\log(1-z/\rho)
 +2\log z\operatorname{Li}_2(z/\rho)
 -2\operatorname{Li}_3(z/\rho).
\end{equation}
Differentiation of the polylogarithms cancels the terms proportional
to $\operatorname{Li}_2$ and $\log(1-z/\rho)$, leaving
\begin{equation}
 \frac{\partial\mathcal J(z,\rho)}{\partial z}
 =\frac{\log^2 z}{z-\rho}.
\end{equation}
For the ratio $R$, the inverse transformation is
$q=(1-R)/(a-bR)$. Consequently
\begin{equation}
 \frac{\dd q}{q}=\frac{\dd R}{R-1}
                   -\frac{\dd R}{R-a/b}.
\end{equation}
This second pole is essential. Applying this identity to
$2\log A\log B=\log^2A+\log^2B-\log^2R$ reduces the
product of logarithms to the same primitive. The remaining mixed
term follows from
\begin{equation}
 \frac{\dd}{\dd q}\left[-\ln q\operatorname{Li}_2(aq)
                  +\operatorname{Li}_3(aq)\right]
 =\frac{\ln q\log(1-aq)}q.
\end{equation}
Combining these operations gives
\begin{align}
 \mathcal F(q)&=2\mathcal J(A,1)+2\mathcal J(B,1)
                  -\mathcal J(R,1)+\mathcal J(R,a/b)
                  +\mathcal D(q),\\
 \mathcal D(q)&=2\ln q\left[\operatorname{Li}_2(aq)
                         +\operatorname{Li}_2(bq)\right]
 -2\left[\operatorname{Li}_3(aq)+\operatorname{Li}_3(bq)\right]
 +\frac13\ln^3q,\\
 G_{r,-1,-1}(X)&=-\frac12\left[\mathcal F(q)-\mathcal F(1)\right].
 \label{eq:comp-root-squared-classical}
\end{align}
Indeed, $\mathcal F'(q)=\ln^2[(1-q+q^2)/q]/q$.
The minus sign follows from $\omega_r=-\dd q/q$, and the
subtraction at $q=1$ enforces the original lower limit $X=0$.

The branches in this formula follow directly from the physical
interval. The arguments of $A$ and $B$ lie respectively in
$[-\pi/3,0]$ and $[0,\pi/3]$, so
$\log R=\log A-\log B$ throughout the interval.
The ratio $R$ follows the lower unit-circle arc from $1$ to
$e^{-2\pi\ii/3}$. Neither this arc nor the rotated arc $R/(a/b)$
crosses the polylogarithm cut $[1,\infty)$ in the open interval.
The other polylogarithm arguments are $A$, $B$, $aq$ and $bq$,
which also avoid that cut. The logarithms in $\mathcal J$ are
continuous on these paths. Thus principal values suffice, with no
numerically chosen branch corrections. Although individual terms
are complex, the derivative of $\mathcal F$ is real, so the endpoint
difference in \eqref{eq:comp-root-squared-classical} is real.
This reduction leaves five nonclassical words in the partial sum.

\paragraph{Classical forms for the five remaining words.}
The same primitive also integrates the mixed logarithms generated by
integration by parts. This gives a common classical representation for
all six words, rather than six unrelated integration formulas.
For $d\in\{a,b\}$ and $h\in\{1,-1\}$, define
\begin{align}
 B_d(q)&=1-dq,\qquad R_d(q)=q/B_d(q),\qquad
 T_{hd}(q)=(1-hq)/B_d(q),\\
 N_d(q)&=-\ln q\operatorname{Li}_2(dq)+\operatorname{Li}_3(dq),\\
 P_{hd}(q)&=\tfrac12\bigl[\mathcal J(1-hq,1)+\mathcal J(B_d,1)
             -\mathcal J(T_{hd},1)+\mathcal J(T_{hd},h/d)\bigr],\\
 H_{dc}(q)&=\tfrac12\bigl[\mathcal J(q,c)+\mathcal J(B_d,1-dc)
             -\mathcal J(R_d,c/(1-dc))+\mathcal J(R_d,-1/d)\bigr].
 \label{eq:comp-root-mixed-functions}
\end{align}
The definition of $N_d$ is also used for $d=1,-1$.
In $H_{dc}$ only $c=1,\ii,-\ii$ is needed.
To derive these expressions, apply the square-of-a-ratio identity
used above to each pair of logarithms. The transformed measures
include both the finite pole and the image of infinity. Direct
differentiation then gives
\begin{align}
 N_d'(q)&=\frac{\ln q\log(1-dq)}q,\\
 P_{hd}'(q)&=\frac{\log(1-hq)\log(1-dq)}q,\\
 H_{dc}'(q)&=\frac{\ln q\log(1-dq)}{q-c}.
 \label{eq:comp-root-mixed-derivatives}
\end{align}
These identities specify the primitives independently of any cross-section
formula from the literature.

Some individual terms have a singular expression at $q=1$, although
the required combinations are finite. Their boundary values must therefore
be evaluated before numerical subtraction. The necessary limits are
\begin{align}
 \mathcal J(0,\rho)&=0\quad(\rho\ne0),\qquad
 \mathcal J(1,1)=-2\zeta(3),\\
 P_{1d}(1)&=\tfrac12\mathcal J(1-d,1),\\
 H_{d1}(1)&=\tfrac12\bigl[-2\zeta(3)
       +\log^2(1-d)\log(-d)+4\zeta(2)\log(1-d)
       +\mathcal J(1/(1-d),-1/d)\bigr].
 \label{eq:comp-root-mixed-boundary}
\end{align}
For example, the two divergent logarithms in $H_{d1}$ occur in
$\mathcal J(B_d,1-d)-\mathcal J(R_d,1/(1-d))$.
As $q\to1^-$, their difference tends to $\log(-d)$.
Their coefficients both tend to $\log^2(1-d)$ because
$\log(1/(1-d))=-\log(1-d)$. The dilogarithms tend to
$\zeta(2)$, giving the displayed $4\zeta(2)\log(1-d)$ term.
This determines the boundary constant analytically.

Put $S(q)=N_a(q)+N_b(q)$ and define
\begin{align}
 U(q)&=-2\sum_{d=a,b}P_{1d}(q)+S(q)+2N_1(q)-\tfrac13\ln^3q,\\
 V(q)&=\ln q\operatorname{Li}_2(-X)
       +2\sum_{d=a,b}H_{d1}(q)-S(q)
       -2\mathcal J(q,1)+\tfrac13\ln^3q,\\
 W(q)&=2\sum_{d=a,b}P_{-1,d}(q)-S(q)-2N_{-1}(q)
       +\tfrac13\ln^3q+\ln4\,E(q),\\
 E(q)&=\operatorname{Li}_2(aq)+\operatorname{Li}_2(bq)
                 +\tfrac12\ln^2q,\\
 Z(q)&=-\ln q\,D(X)
       +\sum_{c=\ii,-\ii}\left[\sum_{d=a,b}H_{dc}(q)-\mathcal J(q,c)\right]
       -S(q)+\tfrac13\ln^3q,\\
 D(X)&=\operatorname{Li}_2(-1-X)+\ln(1+X)\ln(2+X)+\frac{\pi^2}{12}.
\end{align}
All occurrences of $X$ here mean $(1-q)^2/q$.
Write $\Delta f=f(q)-f(1)$, with the finite boundary values just derived.
The five reductions are
\begin{align}
 G_{r,0,-1}(X)&=\Delta V,\\
 G_{r,-1,0}(X)&=\Delta U-\Delta V,\\
 G_{0,r,-1}(X)&=\ln X\,\mathcal G(X)-\Delta U,\\
 G_{-4,r,-1}(X)&=\ln\!\left(\frac{X+4}{4}\right)\mathcal G(X)+\Delta W,\\
 G_{r,-2,-1}(X)&=\Delta Z.
 \label{eq:comp-root-five-classical}
\end{align}
Here $\mathcal G(X)=G_{r,-1}(X)$ is the classical dilogarithm
combination in \eqref{eq:comp-ir-short-function}.

For completeness, the integration steps fixing these formulas can be
checked without manipulating the full rational coefficients of the cross
section. Since $G_{0,-1}=-\operatorname{Li}_2(-X)$, integration by parts gives
\begin{equation}
 G_{r,0,-1}(X)=\ln q\operatorname{Li}_2(-X)
 +\int_1^q\ln t\ln(1+X(t))
       \left(\frac2{t-1}-\frac1t\right)\dd t.
\end{equation}
Inserting $\ln(1+X(t))=\log(1-at)+\log(1-bt)-\ln t$
and using \eqref{eq:comp-root-mixed-derivatives} gives $\Delta V$.
The shuffle identity gives
$G_{-1,0}+G_{0,-1}=\ln X\ln(1+X)$, whose integral with
$\omega_r$ is $\Delta U$. Multiplication of $G_{r,-1}$ by
$\ln X$ then gives the third line of
\eqref{eq:comp-root-five-classical}.
For the fourth line, integrate the outer $-4$ kernel by parts and use
\begin{equation}
 \ln\!\left(\frac{X(t)+4}{4}\right)
       =2\ln(1+t)-\ln t-\ln4.
\end{equation}
Finally, $D(X)=G_{-2,-1}(X)$ satisfies
$\dd D/\dd X=\ln(1+X)/(X+2)$. The factorization
$X(t)+2=(t-\ii)(t+\ii)/t$ gives
\begin{equation}
 \frac{\dd D(X(t))}{\dd t}
 =\ln(1+X(t))\left(\frac1{t-\ii}+\frac1{t+\ii}-\frac1t\right).
\end{equation}
Integrating $-D(X(t))\,\dd t/t$ by parts now gives $\Delta Z$.
Thus all constants follow from the original zero lower limit, and
no GPL remains among these six functions.

The complex arguments in the mixed primitives also have fixed physical
branches. For $d=a$, $B_d$ lies below the real axis, while
$q/B_d$ and $(1\pm q)/B_d$ lie above it. Their arguments stay
between $0$ and $\pi/3$, so each logarithm of a ratio is the
difference of the corresponding logarithms. The $d=b$ expressions
are their conjugate continuations. For $c=\pm\ii$, the arguments
$B_a/(1-ac)$ and $R_a(1-ac)/c$ can cross the real axis, but
only below $1$. This follows by multiplying numerator and denominator
by the conjugate denominator. For instance the possible crossing is
at $q=2-\sqrt3$, where $B_a/(1-a\ii)=2\sqrt3-3<1$ and
$R_a(1+a\ii)/(-\ii)=1-2\sqrt3/3<0$. Thus these paths do not cross the
polylogarithm cut. The sole real argument on that cut is $1+q$
in $\mathcal J(1+q,1)$. There the real boundary value is used:
\begin{equation}
 \mathcal J(z,1)=\ln^2z\ln(z-1)
 +2\ln z\operatorname{Re}\operatorname{Li}_2(z)
 -2\operatorname{Re}\operatorname{Li}_3(z),\qquad z>1.
\end{equation}
Its derivative is $\ln^2z/(z-1)$ and its limit at $z=1$ is
$-2\zeta(3)$. This fixes the continuation without a fitted imaginary
part. Equations \eqref{eq:comp-root-squared-classical} and
\eqref{eq:comp-root-five-classical} complete the classical reduction
of this virtual and UV partial sum. Further simplification of their
combined coefficients and arguments remains a separate step.

\paragraph{Combining the root terms in the cross section.}
The preceding conversion to classical functions does not by itself shorten
the cross section. Their coefficients must also be combined. Let
$C_{\mathrm{six}}(X)$ denote just the contribution of the six words listed
above to $\widehat C_{V+\mathrm{UV}}$. Their original rational coefficients
have the common factor $Y/[X^4(X+4)]$. After inserting their classical
forms, the coefficients of $\Delta U$ and $\Delta V$ are in the exact
ratio $-3:4$. Define $K(q)=4V(q)-3U(q)$ and use the polynomials in
\eqref{eq:comp-root-compact-polynomials}.
The six contributions then combine into
\begin{equation}
 \begin{aligned}
 C_{\mathrm{six}}(X)=\frac{Y}{X^4(X+4)}\bigg\{&
 -\frac32p_F\Delta\mathcal F+p_K\Delta K
 -6p_W\Delta W-8p_Z\Delta Z\\
 &+\left[3p_0\ln X-6p_W\ln\!\left(\frac{X+4}{4}\right)\right]
       \mathcal G(X)\bigg\}.
 \end{aligned}
 \label{eq:comp-root-combined-coefficients}
\end{equation}
Thus four boundary differences suffice for this contribution. In
particular, $U$ and $V$ need not be evaluated separately.
There is also a cancellation inside $K$: the $8\mathcal J(q,1)$
generated by $8\sum_d H_{d1}$ cancels its explicit
$-8\mathcal J(q,1)$ term. This follows by expanding
\eqref{eq:comp-root-mixed-functions}, before evaluating any special
function. Equation~\eqref{eq:comp-root-combined-coefficients} is an exact
coefficient collection of the six original terms in this diagram subset.
The other virtual, soft, and real terms must still be added using
the common-regulator assembly \eqref{eq:comp-assembly-common-regulator}.
